%! TEX root = main.tex  

\begin{abstract}
This project investigates mixed music as an interactive performance system by coupling an acoustic cello with a live electronics environment whose behavior depends on what happens on stage, rather than being fixed in advance. Conceived for two performers (cellist and electronic performer operating a SuperCollider patch), the work treats cello and electronics as an expanded instrument: the electronic layer does not merely accompany, but extends the cello through real-time capture, transformation and re-projection of its sound. Starting from culturally recognizable instrumental fragments, the piece progressively shifts the listener’s attention from score-oriented decoding (style, gesture, expectation) toward a more reduced listening, where timbre, grain, density and the perceived “weight” of sound become primary. The extra-musical image guiding the form is memory: initially faithful (echo/repetition), then selective (reconstruction) and finally unstable (erosion and saturation) until the original identity survives only as a trace. All sonic material derives from the cello and is organized into a palette of gesture families (melodic/tonic, articulated and noisy/inharmonic objects) shaped through morphological criteria such as mass, matter and dynamic evolution. The macroform follows a four-part trajectory (presence, echo/reconstruction, counterpoint with the past, dissolution/coda), realized through buffer-based strategies including delay/accumulation, looping and layering, fragmentation and scratching and a final drift toward granular textures.
\end{abstract}
